\uuid{9UyE}
\exo7id{3614}
\titre{exo7 3614}
\auteur{quercia}
\organisation{exo7}
\datecreate{2010-03-10}
\isIndication{false}
\isCorrection{true}
\chapitre{Réduction d'endomorphisme, polynôme annulateur}
\sousChapitre{Sous-espace stable}
\cpge{mpsi-mp2i.sous-espaces-affines, ptsi.espaces-vectoriels.b, mp-mpi.reduction-endomorphismes.b, pc.reduction.a, psi.reduction-endomorphismes.a, pt.reduction-endomorphismes.a, bcpst2.valeurs-propres-vecteurs-propres.a}
\module{Algèbre}
\niveau{L2}
\difficulte{}

\contenu{
\question{
Soit $E = \R^3$ et $H:x+2y+3z=1$ un plan {\it affine\/} de $E$.
Montrer que si $H$ est stable par $f\in\mathcal{L}(E)$ alors $1$ est
valeur propre de~$f$.
}
\reponse{
Soit $\varphi(x,y,z) = x+2y+3z$. $f$ conserve la surface de
niveau $\varphi=1$ donc par linéarité $f\circ\varphi=\varphi$ et
$\varphi$ est vecteur propre de ${}^tf$.
}
}
